\section{Analysis and Falsification}

\subsection{Matched-context analysis}

Repeated blocks are nested within participants. The primary model estimates condition effects for S, E, and A while accounting for order, session, challenge type, and prespecified participant-level covariates when justified.

Movement exposure can be tested for superiority because the manipulation is designed to increase physical engagement. Cognitive performance should instead use a preregistered non-inferiority framework when the scientific aim is to determine whether embodiment or AR can be added without unacceptable loss of control performance.

Let \(\theta_{ES}\) denote the mean embodied-minus-stationary difference on a cognitive outcome scaled so that positive values indicate better performance. For non-inferiority margin \(\Delta_{NI}>0\), E is non-inferior to S only if the prespecified confidence interval excludes effects worse than \(-\Delta_{NI}\). The margin must be justified before outcome inspection.

The AR-specific cognitive comparison uses the analogous \(A-E\) contrast. A significant \(A-S\) difference alone cannot identify whether the effect is attributable to movement or AR.

\subsection{Recovery analysis}

Monitoring-anchored recovery latency \(T_{\mathrm{rec}}\) remains a time-to-event outcome with right-censoring when the recovery criterion is not reached. Context can enter a mixed survival or accelerated-failure-time model as a within-participant predictor. Because movement or AR can alter pre-monitoring transition time, \(T_{\mathrm{pre}}\) and \(T_{\mathrm{total}}\) are analyzed separately rather than absorbed into \(T_{\mathrm{rec}}\). A context effect is interpreted together with cognitive accuracy, challenge equivalence, tracking quality, and safety burden.

\subsection{Transfer analysis}

A later longitudinal contextual-transfer study requires an independent endpoint that is not itself an AR or trained movement task. The relevant estimand is the between-group difference in change or follow-up performance between varied-context and fixed-context training.

\subsection{Falsification criteria}

The embodied or AR transfer claim is weakened when:

\begin{enumerate}
\item E increases movement but violates cognitive non-inferiority relative to S;
\item A differs from S but not from E, indicating that the effect is explained by embodiment rather than AR;
\item apparent AR benefit disappears after accounting for novelty, order, or technical tracking quality;
\item recovery changes do not generalize to independent transfer outcomes;
\item safety, fatigue, visual burden, or environmental constraints make the context unsuitable for reliable home use.
\end{enumerate}
